\documentclass[12pt]{article}
\usepackage[utf8]{inputenc}
\usepackage{amsmath}
\usepackage{amssymb}
\usepackage{geometry}
\usepackage{graphicx}
\usepackage{hyperref}
\usepackage{physics}
\usepackage{caption}
\usepackage{subcaption}
\usepackage{enumitem}
\usepackage{titlesec}
\usepackage{float}
\usepackage[table]{xcolor}
\usepackage{colortbl}
\usepackage{array}
\usepackage{makecell}
\usepackage{listings}

\lstset{
	numbers=left,
	tabsize=1,
	tab=1,
	xleftmargin=1pt,
	basicstyle=\ttfamily,
	breaklines=true,
	frame=single,
}

\geometry{
	a4paper,
	left=25mm,
	right=25mm,
	top=25mm,
	bottom=25mm
}

\title{
	\textbf{Principles of Data Communications and Networks Modeling (CS-6865)} \\[2mm]
	\text{Lab Report \#3: Investigation of network congestion 1}
}
\author{Abid Hasan Muin (3804634)}
\date{Due date: October 28, 2025}


\titleformat{\section}{\large\bfseries}{Exercise \thesection:}{1em}{}
\setcounter{section}{0}

\begin{document}
	
	\maketitle
	
	% Exercise 1
	\section{Plotting throughput with Xgraph, custom trace files}
	
	% Exercise 1.1
	\subsection{Screenshot of Xgraph Output}
	
	\begin{figure}[H]
		\centering
		\includegraphics[width=\textwidth]{exer-1_1_a.png}
		\caption{Xgraph output}
		\label{fig:exer_3_1_a}
	\end{figure}
	
	% Exercise 1.2
	\subsection{Explanation of the custom trace file}
	
	\begin{lstlisting}
	0 0.0
	0.5 0.0
	1 0.0
	1.5 0.0
	2 0.0
	2.5 0.0
	3 0.0
	3.5 0.0
	4 0.0
	4.5 0.0
	5 0.0
	5.5 0.0
	6 0.0
	6.5 0.0
	7 0.0
	7.5 0.0
	8 0.0
	8.5 0.0
	9 0.0
	9.5 0.0
	10 0.0
	10.5 0.0
	11 0.041599999999999998
	11.5 0.099199999999999997
	12 0.099199999999999997
	12.5 0.099199999999999997
	13 0.1024
	13.5 0.099199999999999997
	14 0.099199999999999997
	14.5 0.099199999999999997
	\end{lstlisting}
	
	\noindent
	\textbf{Explanation:}
	The trace file \textit{out0.tr} contains two columns: the first column is the simulation time in seconds, and the second column is the measured throughput at that time. For example, at 11 seconds, the throughput is 0.0416 Mbps, and at 13 seconds, it is 0.1024 Mbps. Initially, the throughput is 0 because the flow has not started. As the simulation progresses, the throughput increases and fluctuates slightly due to packet transmission intervals.
	
	% Exercise 1.3
	\subsection{Add NAM trace}
	
	\begin{figure}[H]
		\centering
		\includegraphics[width=\textwidth]{exer-1_1_c.png}
		\caption{NAM trace}
		\label{fig:exer_3_1_c}
	\end{figure}
	
	\textbf{Modified code:}
	\begin{lstlisting}
	#Create a simulator object
	set ns [new Simulator]
	
	# 1.c)
	# Create a NAM trace file
	set nf [open out.nam w]
	$ns namtrace-all $nf
	# 1.c)
	
	#Open the output files
	set f0 [open out0.tr w]
	set f1 [open out1.tr w]
	set f2 [open out2.tr w]
	
	#Create 5 nodes
	set n0 [$ns node]
	set n1 [$ns node]
	set n2 [$ns node]
	set n3 [$ns node]
	set n4 [$ns node]
	
	#Connect the nodes
	$ns duplex-link $n0 $n3 1Mb 100ms DropTail
	$ns duplex-link $n1 $n3 1Mb 100ms DropTail
	$ns duplex-link $n2 $n3 1Mb 100ms DropTail
	$ns duplex-link $n3 $n4 1Mb 100ms DropTail
	
	#Define a 'finish' procedure
	proc finish {} {
		global f0 f1 f2 ns nf
		#Close the Xgraph output files
		close $f0
		close $f1
		close $f2
		
		$ns flush-trace
		# Close the NAM file
		close $nf
		#Call xgraph to display the results
		exec xgraph out0.tr out1.tr out2.tr -geometry 800x400 &
		# Launch NAM for animation
		exec nam out.nam &
		exit 0
	}
	
	
	#Define a procedure that attaches a UDP agent to a previously created node
	#'node' and attaches an Expoo traffic generator to the agent with the
	#characteristic values 'size' for packet size 'burst' for burst time,
	#'idle' for idle time and 'rate' for burst peak rate. The procedure connects
	#the source with the previously defined traffic sink 'sink' and returns the
	#source object.
	proc attach-expoo-traffic { node sink size burst idle rate } {
		#Get an instance of the simulator
		set ns [Simulator instance]
		
		#Create a UDP agent and attach it to the node
		set source [new Agent/UDP]
		$ns attach-agent $node $source
		
		#Create an Expoo traffic agent and set its configuration parameters
		set traffic [new Application/Traffic/Exponential]
		$traffic set packetSize_ $size
		$traffic set burst_time_ $burst
		$traffic set idle_time_ $idle
		$traffic set rate_ $rate
		
		# Attach traffic source to the traffic generator
		$traffic attach-agent $source
		#Connect the source and the sink
		$ns connect $source $sink
		return $traffic
	}
	
	
	#Define a procedure which periodically records the bandwidth received by the
	#three traffic sinks sink0/1/2 and writes it to the three files f0/1/2.
	proc record {} {
		global sink0 sink1 sink2 f0 f1 f2
		#Get an instance of the simulator
		set ns [Simulator instance]
		#Set the time after which the procedure should be called again
		set time 0.5
		#How many bytes have been received by the traffic sinks?
		set bw0 [$sink0 set bytes_]
		set bw1 [$sink1 set bytes_]
		set bw2 [$sink2 set bytes_]
		#Get the current time
		set now [$ns now]
		#Calculate the bandwidth (in MBit/s) and write it to the files
		puts $f0 "$now [expr $bw0/$time*8/1000000]"
		puts $f1 "$now [expr $bw1/$time*8/1000000]"
		puts $f2 "$now [expr $bw2/$time*8/1000000]"
		#Reset the bytes_ values on the traffic sinks
		$sink0 set bytes_ 0
		$sink1 set bytes_ 0
		$sink2 set bytes_ 0
		#Re-schedule the procedure
		$ns at [expr $now+$time] "record"
	}
	
	
	#Create three traffic sinks and attach them to the node n4
	set sink0 [new Agent/LossMonitor]
	set sink1 [new Agent/LossMonitor]
	set sink2 [new Agent/LossMonitor]
	$ns attach-agent $n4 $sink0
	$ns attach-agent $n4 $sink1
	$ns attach-agent $n4 $sink2
	
	#Create three traffic sources
	set source0 [attach-expoo-traffic $n0 $sink0 200 2s 1s 100k]
	set source1 [attach-expoo-traffic $n1 $sink1 200 2s 1s 200k]
	set source2 [attach-expoo-traffic $n2 $sink2 200 2s 1s 300k]
	
	#Start logging the received bandwidth
	$ns at 0.0 "record"
	#Start the traffic sources
	$ns at 10.0 "$source0 start"
	$ns at 10.0 "$source1 start"
	$ns at 10.0 "$source2 start"
	#Stop the traffic sources
	$ns at 50.0 "$source0 stop"
	$ns at 50.0 "$source1 stop"
	$ns at 50.0 "$source2 stop"
	#Call the finish procedure after 60 seconds simulation time
	$ns at 60.0 "finish"
	
	#Run the simulation
	$ns run
	\end{lstlisting}
	
	% Exercise 2
	\section{Two Competing CBR Flows}
	
	% Exercise 2.1
	\subsection{Two Screenshots of Simulation}
	
	\textbf{Calculation for Interval value:}	\\
	Packet size = 500 bytes = 500 $\times$ 8 = 4000 bits \\
	Rate = 800 Kbps = 800 $\times$ 1000 = 800,000 bps \\
	$\text{Interval} = \frac{\text{Packet size}}{\text{Rate}} = \frac{4000}{800000} = 0.005~\text{seconds}$
	
	\begin{figure}[H]
		\centering
		\includegraphics[width=\textwidth]{exer-2_a_1.png}
		\caption{Single CBR Flow}
		\label{fig:exer-2_a_1}
	\end{figure}
	
	\begin{figure}[H]
		\centering
		\includegraphics[width=\textwidth]{exer-2_a_2.png}
		\caption{Double CBR Flow}
		\label{fig:exer-2_a_2}
	\end{figure}
	
	% Exercise 2.2
	\subsection{Throughput Analysis}
	
	\textbf{Modified code:}
	\begin{lstlisting}
	# =========================
	# Exercise 2(b): Two Competing CBR Flows - Throughput Measurement
	# =========================
	
	# Create simulator
	set ns [new Simulator]
	
	# Define colors for flows (this affects NAM visualization)
	$ns color 0 Red
	$ns color 1 Blue
	
	# Create NAM trace file
	set nf [open out.nam w]
	$ns namtrace-all $nf
	
	# Create Xgraph output files
	set f0 [open out0.tr w]
	set f1 [open out1.tr w]
	
	# =========================
	# Create nodes
	# =========================
	set n0 [$ns node]
	set n1 [$ns node]
	set n2 [$ns node]
	set n3 [$ns node]
	set n4 [$ns node]
	set n5 [$ns node]
	
	# =========================
	# Connect nodes (1Mbps, 10ms, DropTail)
	# =========================
	$ns duplex-link $n0 $n2 1Mb 10ms DropTail
	$ns duplex-link $n1 $n2 1Mb 10ms DropTail
	$ns duplex-link $n2 $n3 1Mb 10ms DropTail
	$ns duplex-link $n3 $n4 1Mb 10ms DropTail
	$ns duplex-link $n3 $n5 1Mb 10ms DropTail
	
	# Orient links for better visualization
	$ns duplex-link-op $n0 $n2 orient right-down
	$ns duplex-link-op $n1 $n2 orient right-up
	$ns duplex-link-op $n2 $n3 orient right
	$ns duplex-link-op $n3 $n4 orient right-up
	$ns duplex-link-op $n3 $n5 orient right-down
	
	# =========================
	# Create traffic sinks
	# =========================
	set sink0 [new Agent/LossMonitor]
	$ns attach-agent $n4 $sink0
	
	set sink1 [new Agent/LossMonitor]
	$ns attach-agent $n5 $sink1
	
	# =========================
	# Create UDP agents and CBR sources
	# =========================
	# Flow 0: n0 -> n4
	set udp0 [new Agent/UDP]
	$udp0 set class_ 0
	$ns attach-agent $n0 $udp0
	set cbr0 [new Application/Traffic/CBR]
	$cbr0 set packetSize_ 500
	$cbr0 set interval_ 0.005     ;# 800 Kbps
	$cbr0 attach-agent $udp0
	$ns connect $udp0 $sink0
	
	# Flow 1: n1 -> n5
	set udp1 [new Agent/UDP]
	$udp0 set class_ 1
	$ns attach-agent $n1 $udp1
	set cbr1 [new Application/Traffic/CBR]
	$cbr1 set packetSize_ 500
	$cbr1 set interval_ 0.005     ;# 800 Kbps
	$cbr1 attach-agent $udp1
	$ns connect $udp1 $sink1
	
	# =========================
	# Record throughput for Xgraph
	# =========================
	proc record {} {
		global ns sink0 sink1 f0 f1
		set time 0.5
		set now [$ns now]
		
		set bw0 [$sink0 set bytes_]
		set bw1 [$sink1 set bytes_]
		puts $f0 "$now [expr $bw0/$time*8/1000000]"
		puts $f1 "$now [expr $bw1/$time*8/1000000]"
		
		$sink0 set bytes_ 0
		$sink1 set bytes_ 0
		$ns at [expr $now+$time] "record"
	}
	
	$ns at 0.0 "record"
	
	# =========================
	# Flow start/stop times
	# =========================
	$ns at 1.0 "$cbr0 start"
	$ns at 20.0 "$cbr0 stop"
	
	$ns at 5.0 "$cbr1 start"
	$ns at 15.0 "$cbr1 stop"
	
	# =========================
	# Finish routine
	# =========================
	proc finish {} {
		global f0 f1 nf
		close $f0
		close $f1
		close $nf
		
		exec xgraph out0.tr out1.tr -geometry 1920x1080 &
		exec nam out.nam &
		exit 0
	}
	
	# Call finish at end of simulation
	$ns at 25.0 "finish"
	
	# =========================
	# Run simulation
	# =========================
	$ns run
	\end{lstlisting}
	
	\begin{figure}[H]
		\centering
		\includegraphics[width=\textwidth]{exer-2_b_1.png}
		\caption{Xgraph output}
		\label{fig:exer-2_b_1}
	\end{figure}
	
	\noindent
	\textbf{Comment:}	\\
	When only flow n0–n4 was active (1 s – 5 s), the measured throughput was about 0.8 Mbps with no packet loss. When flow n1–n5 started at 5 s, both CBR sources jointly offered 1.6 Mbps to the 1 Mbps bottleneck link, causing severe congestion. The earlier flow (n0–n4) maintained nearly full throughput, while the later flow (n1–n5) experienced heavy losses. The total delivered throughput of $\approx$ 1 Mbps matches the bottleneck capacity. This illustrates how UDP traffic without congestion control can create unfair sharing and packet drops on shared links.
	
	% Exercise 2.3
	\subsection{Packet Drop Analysis}
	
	\textbf{Flow 0 (started first):}	\\
	Received 3805 packets	\\
	Average throughput: $\approx$0.8 Mbps overall, $\approx$0.64 Mbps during competition	\\
	Experienced: 0 drops	\\
	
	\noindent
	\textbf{Flow 1 (started later):}	\\
	Received only 523 packets	\\
	Average throughput: ~0.21 Mbps during competition	\\
	Experienced: 1492 drops (100\% of all drops)	\\
	
	\noindent
	\textbf{The Unfairness is Even Worse Than Expected:}	\\
	Not only does Flow 1 suffer all the drops, but it also:	\\
	Gets only 25\% of the bandwidth it should get	\\
	While Flow 0 maintains 80\% of its target bandwidth	\\
	3:1 bandwidth ratio instead of the expected 1:1	\\
	
	\noindent
	\textbf{Why This Happens:}	\\
	The DropTail queue + timing synchronization creates a vicious cycle:	\\
	Flow 0's packets arrive first $\rightarrow$ get queued	\\
	Flow 1's packets arrive second $\rightarrow$ queue full $\rightarrow$ dropped	\\
	Flow 0 packets get transmitted $\rightarrow$ make space in queue	\\
	Next round: Flow 0 packets again arrive first $\rightarrow$ repeat	\\
	Flow 1 can only squeeze in packets during rare moments when its timing happens to beat Flow 0, which is why it gets ~500 packets instead of 0, but nowhere near the ~1600 it should get.
	
	
	% Exercise 3
	\section{CBR flow competing with FTP over TCP flow}
	
	% Exercise 3.1
	\subsection{Two Screenshots of Simulation}
	
	\begin{figure}[H]
		\centering
		\includegraphics[width=\textwidth]{exer-3_a_1.png}
		\caption{Single Flow}
		\label{fig:exer-3_a_1}
	\end{figure}
	
	\begin{figure}[H]
		\centering
		\includegraphics[width=\textwidth]{exer-3_a_2.png}
		\caption{Double Flow}
		\label{fig:exer-3_a_2}
	\end{figure}
	
	% Exercise 3.2
	\subsection{Throughput Analysis}
	
	\textbf{Modified code:}
	\begin{lstlisting}
	# Exercise 3(b): TCP/FTP competing with CBR/UDP
	
	# Create simulator
	set ns [new Simulator]
	
	# Define colors for flows (this affects NAM visualization)
	$ns color 0 Blue;   # TCP
	$ns color 1 Red;    # CBR
	
	# Create NAM trace file
	set nf [open out.nam w]
	$ns namtrace-all $nf
	
	# Create trace file for analysis
	set tracefile [open out.tr w]
	$ns trace-all $tracefile
	
	# Create Xgraph output files
	set f0 [open out0.tr w];    # TCP
	set f1 [open out1.tr w];    # CBR
	
	# Create nodes
	set n0 [$ns node]
	set n1 [$ns node]
	set n2 [$ns node]
	set n3 [$ns node]
	set n4 [$ns node]
	set n5 [$ns node]
	
	# Connect nodes (1Mbps, 10ms, DropTail)
	$ns duplex-link $n0 $n2 1Mb 10ms DropTail
	$ns duplex-link $n1 $n2 1Mb 10ms DropTail
	$ns duplex-link $n2 $n3 1Mb 10ms DropTail
	$ns duplex-link $n3 $n4 1Mb 10ms DropTail
	$ns duplex-link $n3 $n5 1Mb 10ms DropTail
	
	# Orient links for better visualization
	$ns duplex-link-op $n0 $n2 orient right-down
	$ns duplex-link-op $n1 $n2 orient right-up
	$ns duplex-link-op $n2 $n3 orient right
	$ns duplex-link-op $n3 $n4 orient right-up
	$ns duplex-link-op $n3 $n5 orient right-down
	
	
	# Flow 0: TCP / FTP (n0 -> n4)
	set tcp0 [new Agent/TCP]
	$tcp0 set class_ 0
	$tcp0 set window_ 100
	$tcp0 set packetSize_ 500
	$ns attach-agent $n0 $tcp0
	
	set sink0 [new Agent/TCPSink]
	$ns attach-agent $n4 $sink0
	
	$ns connect $tcp0 $sink0
	
	set ftp0 [new Application/FTP]
	$ftp0 attach-agent $tcp0
	
	
	# Flow 1: CBR / UDP (n1 -> n5)
	set udp1 [new Agent/UDP]
	$udp1 set class_ 1
	$ns attach-agent $n1 $udp1
	
	set sink1 [new Agent/LossMonitor]
	$ns attach-agent $n5 $sink1
	
	$ns connect $udp1 $sink1
	
	set cbr1 [new Application/Traffic/CBR]
	$cbr1 set packetSize_ 500
	$cbr1 set interval_ 0.005     ;# 800 Kbps
	$cbr1 attach-agent $udp1
	
	
	# Record throughput for Xgraph
	proc record {} {
		global ns f0 f1 sink0 sink1
		
		set timeInterval 0.5
		set now [$ns now]
		
		# Calculate throughput (Mbps)
		set bw0 [expr [$sink0 set bytes_]*8.0/($timeInterval*1000000)]
		set bw1 [expr [$sink1 set bytes_]*8.0/($timeInterval*1000000)]
		
		puts $f0 "$now $bw0"
		puts $f1 "$now $bw1"
		
		# Reset counters for next interval
		$sink0 set bytes_ 0
		$sink1 set bytes_ 0
		
		$ns at [expr $now+$timeInterval] "record"
	}
	
	
	$ns at 0.0 "record"
	
	# Flow start/stop times
	# n0-n4
	$ns at 1.0 "$ftp0 start"
	$ns at 20.0 "$ftp0 stop"
	# n1-n5
	$ns at 5.0 "$cbr1 start"
	$ns at 15.0 "$cbr1 stop"
	
	# Finish routine
	proc finish {} {
		global ns nf f0 f1 tracefile
		
		$ns flush-trace
		close $tracefile
		close $f0
		close $f1
		close $nf
		
		exec xgraph out0.tr out1.tr -geometry 1920x1080 -x "Time (s)" -y "Throughput (Mbps)" &
		exec nam out.nam &
		exit 0
	}
	
	# Call finish at end of simulation
	$ns at 25.0 "finish"
	
	# Run simulation
	$ns run
	\end{lstlisting}
	
	
	\subsubsection{Throughput Analysis Results}
	
	\begin{figure}[H]
		\centering
		\includegraphics[width=\textwidth]{exer-3_b_1.png}
		\caption{Xgraph screenshot}
		\label{fig:exer-3_b_1}
	\end{figure}
	
	\textbf{Single Flow Transmission Periods:}	\\
	\noindent	\\
	Period 1.0-5.0 seconds (Only TCP/FTP active):	\\
	TCP Throughput: $\approx$0.95-1.0 Mbps	\\
	Observation: TCP achieves nearly 100\% utilization of the 1 Mbps bottleneck link when operating alone, demonstrating efficient congestion control and window management.\\
	\noindent	\\
	Period 15.0-20.0 seconds (Only TCP/FTP active):	\\
	TCP Throughput: ~0.95-1.0 Mbps	\\
	Observation: TCP quickly recovers full link utilization after the CBR flow terminates at 15 seconds.	\\
	\noindent	\\
	Competing Flows Period (5.0-15.0 seconds):	\\
	Individual Flow Throughput:	\\
	CBR/UDP Flow: Maintains constant $\approx$0.8 Mbps (as configured)	\\
	TCP/FTP Flow: Reduced $\approx$0.2 Mbps	\\
	Observation: The CBR flow dominates due to its aggressive, non-responsive nature	\\
	\noindent	\\
	Cumulative Throughput:	\\
	Total: $\approx$1.0 Mbps	\\
	Observation: The bottleneck link remains fully utilized, demonstrating efficient resource sharing despite the competition.	\\	\\
	\noindent	\\
	\textbf{Key Observations from Throughput Graphs}	\\
	TCP Adaptive Behavior:
	\begin{itemize}
		\item Shows significant throughput variation and reduction when competing with CBR
		\item Demonstrates congestion avoidance by backing off transmission rate
		\item Quickly recovers when competition ceases
	\end{itemize}
	
	\noindent
	CBR Non-Responsive Nature:
	\begin{itemize}
		\item Maintains constant 0.8 Mbps rate regardless of network conditions
		\item Does not adapt to congestion, causing unfair bandwidth sharing
	\end{itemize}
	
	\noindent
	Link Utilization Efficiency:
	\begin{itemize}
		\item The bottleneck link maintains $\approx$100\% utilization throughout simulation
		\item No significant bandwidth wastage observed
	\end{itemize}
	
	\noindent
	Fairness Issues:
	\begin{itemize}
		\item CBR/UDP consumes 80\% of available bandwidth while TCP adapts to use only 20\%
		\item This demonstrates the inherent unfairness when responsive TCP competes with non-responsive UDP traffic.
	\end{itemize}
	
	\noindent
	The throughput graphs clearly illustrate the fundamental differences between congestion-responsive (TCP) and non-responsive (CBR/UDP) traffic, highlighting important considerations for network protocol design and quality of service management.
	
	% Exercise 3.3
	\subsection{Packet Drop Analysis}
	
	To analyze packet drops from the trace file, I used the following approach:
	
	\begin{lstlisting}[language=bash]
	# Count total dropped packets
	grep "\^d" out.tr | wc -l
	
	# Count TCP dropped packets
	grep "\^d.*tcp" out.tr | wc -l
	
	# Count CBR dropped packets  
	grep "\^d.*cbr" out.tr | wc -l
	\end{lstlisting}
	
	\noindent
	\textbf{Packet Drop Results:}
	
	\noindent
	Total Dropped Packets:	\\
	Overall 106 packets were dropped during the simulation.	\\
	
	\noindent
	Flow-Specific Packet Drops:	\\
	TCP/FTP Flow (n0-n4): 60 packets dropped (56.6\% of total drops)	\\
	CBR/UDP Flow (n1-n5): 46 packets dropped (43.4\% of total drops)	\\
	
	\noindent
	\textbf{Analysis and Comments:}	\\
	During the period when only the TCP/FTP flow was active (1 s – 5 s), no packet drops occurred since the link bandwidth was sufficient. When the UDP/CBR flow started at 5 s, congestion built up at the bottleneck link (n2–n3). The queue overflowed, resulting in a total of 106 packet drops. About 56.6\% of the dropped packets belonged to the TCP flow, while 43.4\% of UDP packets were lost. This shows that the non-adaptive UDP source consumed most of the bandwidth and forced TCP to back off.
	
	% Exercise 3.4
	\subsection{Comment on Exercise 2 and Exercise 3 Simulation Result}
	In Exercise 2, both UDP/CBR flows transmitted at fixed rates exceeding the bottleneck capacity, leading to persistent congestion and numerous packet drops for both flows. In contrast, Exercise 3 introduced a TCP flow competing with a CBR flow. Because TCP implements congestion control, it reduced its sending rate whenever losses occurred, while the UDP/CBR source continued transmitting at a constant rate. As a result, the total number of dropped packets decreased in Exercise 3 compared with Exercise 2. However, most of the remaining drops were experienced by the TCP flow, which yielded bandwidth to the unresponsive UDP source. The cumulative throughput during the overlap period remained close to the link capacity $\approx$1 Mbps, confirming full utilization but poor fairness.
	
	
	% References
	\begin{thebibliography}{9}
		\bibitem{mark} Mark Greis, Tutorial for the Network Simulator "ns": \url{http://www.isi.edu/nsnam/ns/tutorial/index.html}
		\bibitem{awk} GNU AWK User's Guide: 
		\url{https://www.gnu.org/software/gawk/manual/gawk.html}
	\end{thebibliography}
\end{document}